\section{Discussion}
\label{sec:discussion}

\paragraph{What can and cannot be claimed about the refinement paradox.}
After the audit, the strongest empirical statement supported by
the data is: \emph{In short-answer extractive QA on SQuAD with
Mistral-7B, aggressive per-passage refinement produces a
faithfulness change whose magnitude depends substantially on the
choice of faithfulness scorer.} Under DeBERTa-v3 the magnitude is
small ($0.011$), the per-seed contrasts are mostly non-significant,
and a matched-similarity intervention finds no causal role for CCS.
Under a RAGAS-style LLM-judge the magnitude is large ($0.140$).
The two scorers agree at Pearson $r = 0.18$ on the same outputs.
We do not believe a confident claim about the
\emph{magnitude} of the paradox is supported by the data;
we believe a claim about \emph{measurement instability} is.

\paragraph{Why this matters beyond one phenomenon.}
The single-scorer / small-sample / in-distribution-tuning pattern
that produced the published refinement-paradox magnitude is the
modal evaluation pattern across recent RAG-faithfulness papers.
If the same audit were applied to the dozen or so recently-reported
``faithfulness-improving retrieval interventions,'' the magnitudes
would likely be similarly metric-dependent. We propose
pre-registration, multi-scorer reporting, and explicit
cross-dataset $\tau$ disclosure as the minimum bar for future
RAG-faithfulness empirical claims.

\paragraph{The surviving structural finding.}
\S\ref{sec:coherence_noise} preserves a non-trivial result:
coherence-preserving uninformative noise produces a smaller
faithfulness drop than random off-topic noise at matched per-
passage similarity. This says coherence (in some sense, not
necessarily CCS) is real and structurally separable from
similarity. Future work should look for an
operationalization of coherence that survives the matched-similarity
test of \S\ref{sec:causal_intervention} \emph{and} explains the
slope difference of \S\ref{sec:coherence_noise}. Candidate
operationalizations include attention-weighted retrieved-mass
\citep{geva2023attention}, generator-perplexity over retrieval-set
prefixes, or coherence scored by a separate LLM judge with explicit
schema.

\paragraph{What HCPC-v2 actually is, after the audit.}
We retain HCPC-v2 in the released benchmark as a \emph{candidate
deployment policy}, not as a controlled coherence-causality probe.
Empirically HCPC-v2 is competitive on SQuAD (within 0.0016 of
RAPTOR-2L, slightly above CRAG by 0.010) and dominated on HotpotQA
(below CRAG by 0.011 on faithfulness, by 2.5 pp on hallucination,
and by 18\,\% on latency). Whether it is worth deploying is a
function of the deployment regime; the released LangChain
integration makes the policy switchable per query.

\paragraph{Self-RAG's underperformance.}
The harness fairness question raised in \S\ref{sec:headtohead}
deserves emphasis. Self-RAG was trained on a different distribution
than ours and emits reflection tokens that an entailment scorer
treats as additional answer text; the apparent
$0.57$ faithfulness vs.\ $0.71$ for the other methods may
substantially be a harness mismatch rather than a real capability
gap. We report Self-RAG numbers because they are the
matched-harness numbers a downstream practitioner would actually
observe if they swapped Self-RAG into a CRAG-style pipeline,
which is a useful piece of information even if it is not the
fairest possible Self-RAG comparison.

\paragraph{What the audit costs.}
The full audit ran on free hardware: an Apple M4 Air for Ollama-
backed runs and a Kaggle T4\(\times\)2 session for the 8-bit Self-RAG
comparison. Total runtime is well under 25 hours of free Kaggle GPU
plus 10--15 hours of M4 wall-clock, and zero paid API tokens. The
audit pattern is therefore reproducible by any researcher with
free Kaggle access and a workstation, which we view as a
prerequisite for the recommendation to make multi-scorer
pre-registration routine.
